\section{The ordinary theory on a fixed component}
\label{fixed:section}

The operators of Section~\ref{spec:section} were constructed from the
equivariant theory of $W$.  We now identify what their equations mean near
$y=0$.  On a fixed component, quantum Riemann--Roch removes the inverse
Euler twist.  The same transformation turns the logarithm of the shift
operator into a translation generator in $\lambda$.  Its derivative term
cancels the derivative term of the equivariant grading, leaving the
ordinary grading and a scalar multiple of $z^{-1}$.

\begin{proposition}\label{fixed:equivalence}
Fix a germ at $\lambda\ne0$.  Near $Q=y=0$, use the full Novikov
completion with a formal $\log y$ adjoined, allowing polynomial
dependence on $\log y$ in each coefficient.
Let $F=B$ denote the branch near the eigenvalue zero, or let $F=X$ denote
the whole cluster near the eigenvalue $\lambda$.  Use the projectors and
logarithms near $y=0$ of Proposition~\ref{spec:zero}, and let $A_{k,F}$ be
their modes.  Then, as identities over this coefficient ring,
\[
 \mathcal A_W\text{ satisfies every }A_{k,F}\text{ projectively}
 \quad\Longleftrightarrow\quad
 Z_F\text{ satisfies every }\ell_{k,F}\text{ projectively},
 \qquad k\geq-1.
\]
The operators on the left are infinitesimal symplectic for the equivariant
pairing of $W$.  On the right, $Z_F$ is the ordinary descendant potential
with its full cohomology and actual integral curve labels.
\end{proposition}

We first prove the identity of classical operators.  We then explain its
quantization; this second step requires separating the rank factor of the
Euler class before applying quantum Riemann--Roch.

\subsection{Removing the normal bundle from the grading}

Fix $F$ and abbreviate $n=n_F$, $w=w_F$, and $j=j_F$.  Put
\[
 a=w\lambda,\qquad c_N=c_1(N_F)_{\mathrm{ord}},\qquad W_F=nw.
\]
After restriction to the descendant space of $F$, the grading and shift
from Section~\ref{shift:section} are
\begin{align}
 G_F&=z\partial_z+\lambda\partial_\lambda+
        \frac12+\mu_F-\frac n2+
        \frac{R_F+c_N\cup+nw\lambda}{z},
       \label{fixed:restricted-grading}\\
 T_F&=y^j\prod_x
       \begin{cases}
        x+\lambda,&w=1,\\
        (x-\lambda-z)^{-1},&w=-1
       \end{cases}
       e^{z\partial_\lambda}.
       \label{fixed:restricted-shift}
\end{align}
Here $x$ runs over the ordinary Chern roots of $N_F$, and multiplication
by a class is understood in the second formula.  Symmetric expressions
in the roots are interpreted by the splitting principle.  All expansions
in $x$ are finite after evaluation in the cohomology of $F$.

Choose a germ of $\log a$ at a nonzero value of $\lambda$.  Define a
multiplier $\Delta_F$ by
\begin{equation}\label{fixed:Delta}
 \log\Delta_F=
 \sum_x\left\{
   \frac12\log(a+x)
   +\frac1z\int_0^x-\log(a+t)\,dt
   +\sum_{m\geq1}\frac{B_{2m}}{(2m)!}z^{2m-1}
       \partial_x^{2m-1}\bigl(-\log(a+x)\bigr)
          \right\}.
\end{equation}
The $B_{2m}$ are the Bernoulli numbers.  The coefficient of $z^{-1}$ has
positive ordinary cohomological degree, hence is nilpotent as a
multiplication operator.  Both $\Delta_F$ and its inverse consequently
have a finite lower bound on their powers of $z$.  Their positive tails
are interpreted as formal series at $z=0$.

The multiplier maps the ordinary paired loop space of $F$ to the space
with pairing $\eta_F^t$.  Indeed, all terms in
\eqref{fixed:Delta} except the half logarithm are odd in $z$, so that
\[
 \Delta_F(-z)\Delta_F(z)=\prod_x(a+x)=e_T(N_F).
\]
The inverse Euler factor in $\eta_F^t$ therefore cancels this product.
This also fixes the direction of the square-root factor in
\eqref{fixed:Delta}.

For the scalar part of the shift calculation, put
\begin{equation}\label{fixed:primitive}
 L(v)=v-v\log v,\qquad
 \Theta_F(\lambda)=nL(w\lambda).
\end{equation}

\begin{lemma}\label{fixed:classical}
The following identities hold in the formal Laurent calculus of
Lemma~\ref{spec:gauge}:
\begin{align}
 \Delta_F^{-1}G_F\Delta_F
  &=\ell_{0,F}+\lambda\partial_\lambda+\frac{nw\lambda}{z},
       \label{fixed:grading-conjugation}\\
 \Delta_F^{-1}T_F\Delta_F
  &=y^j\exp\left(
       \frac{\Theta_F(\lambda)-\Theta_F(\lambda+z)}{z}
                  \right)e^{z\partial_\lambda}.
       \label{fixed:shift-conjugation}
\end{align}
Under the same change of space, the logarithm prescribed near $y=0$
becomes
\begin{equation}\label{fixed:log-conjugation}
 z\partial_\lambda+j\log y-\Theta_F'(\lambda)
   =z\partial_\lambda+j\log y+nw\log(w\lambda),
\end{equation}
up to a scalar constant determined by the logarithm branch.
\end{lemma}

\begin{proof}
The cohomological part of the commutator with $\mu_F$ differentiates a
function of a Chern root by $x\partial_x$, because $x$ has Hodge type
$(1,1)$.  Apply
\[
 z\partial_z+\lambda\partial_\lambda+x\partial_x
\]
to one summand in \eqref{fixed:Delta}.  The half logarithm contributes
$1/2$.  If $I(a,x)=\int_0^x-\log(a+t)\,dt$, then
$(\lambda\partial_\lambda+x\partial_x)I=I-x$; thus its contribution is
$-x/z$.  Each Bernoulli term has total degree zero.  Summing over the
roots gives $n/2-c_N/z$, which cancels the two corresponding terms in
\eqref{fixed:restricted-grading}.  This proves
\eqref{fixed:grading-conjugation}.

For the shift, add $L(a)/z$ to the summand indexed by $x$ in
\eqref{fixed:Delta}.  The result depends only on $v=a+x$ and equals
\[
 \Phi(v)=\left((z\partial_v)^{-1}-\frac12+
          \sum_{m\geq1}\frac{B_{2m}}{(2m)!}
                         (z\partial_v)^{2m-1}\right)(-\log v),
\]
where the antiderivative in the first term is $L(v)/z$.  The generating
series for the Bernoulli numbers gives
\begin{equation}\label{fixed:Bernoulli-difference}
 \Phi(v+z)-\Phi(v)=-\log v.
\end{equation}
When $w=1$, this difference cancels the factor $v$ in
\eqref{fixed:restricted-shift}.  When $w=-1$, use instead
$\Phi(v-z)-\Phi(v)=\log(v-z)$, which cancels the factor $(v-z)^{-1}$.
The terms $L(a)/z$ that were added account for the exponential in
\eqref{fixed:shift-conjugation}.

It remains to identify the functional logarithm, rather than just one
operator whose exponential has the required form.  Set
\[
 H=z\partial_\lambda-\Theta_F'(\lambda).
\]
Solving its evolution equation gives
\begin{equation}\label{fixed:log-flow}
 e^{tH}=
  \exp\left(
     \frac{\Theta_F(\lambda)-\Theta_F(\lambda+tz)}{z}
       \right)e^{tz\partial_\lambda}.
\end{equation}
For example, differentiating the right side in $t$ gives $H$ times that
side and its value at $t=0$ is the identity.  This is also the exponential
defined by the holomorphic functional calculus of
Lemma~\ref{spec:calculus}: coefficientwise differentiation of its contour
formula gives the same evolution equation.  The constant term of $H$ in
$z$ is the scalar $-\Theta_F'(\lambda)$.  Locally choose a logarithm
which inverts the exponential at that scalar.  The composition rule in
Lemma~\ref{spec:calculus} then identifies $\log(e^H)$ with $H$.

Conjugation of the resolvents by $\Delta_F$ is allowed by
Lemma~\ref{spec:gauge}.  For $F=B$, apply this argument to the divided
shift $y^{-1}T_F$ and then restore $\log y$; for $F=X$, apply it directly
to $T_F$.  These are exactly the logarithm prescriptions of
Proposition~\ref{spec:zero}.  This proves
\eqref{fixed:log-conjugation}, with the stated freedom in its scalar
constant.
\end{proof}

Subtracting $\lambda/z$ times \eqref{fixed:log-conjugation} from
\eqref{fixed:grading-conjugation} now removes
$\lambda\partial_\lambda$.  Thus the operator $D_F$ on its fixed block,
after the descendant change of space and then $\Delta_F$, is
\begin{equation}\label{fixed:ordinary-grading}
 \ell_{0,F}+\frac{b_F}{z},\qquad
 b_F=\lambda\bigl(nw-j\log y-nw\log(w\lambda)\bigr),
\end{equation}
up to a scalar multiple of $\lambda/z$ from a different logarithm
choice.  In particular, the remaining operator differentiates none of
$\lambda$, $y$, or the Novikov variables.

For a scalar $b$ independent of $z$, write
\[
 \ell_{-1,F}^{(b)}=z^{-1},\qquad
 \ell_{k,F}^{(b)}=z^{-1}(z\ell_{0,F}+b)^{k+1}\quad(k\geq0).
\]
The binomial identity
\begin{equation}\label{fixed:binomial}
 \ell_{k,F}^{(b)}=
   \sum_{m=-1}^{k}\binom{k+1}{m+1}b^{k-m}\ell_{m,F}
\end{equation}
is triangular with diagonal entries one.  Hence these operators are
infinitesimal symplectic, and their simultaneous projective equations
are equivalent to those of the ordinary modes.  We have proved this
identification at the classical level.  To obtain the corresponding
statement for the potentials, we next implement it by quantum
Riemann--Roch.

\subsection{Quantizing the comparison}

The logarithm in \eqref{fixed:Delta} contains
$-\log(w\lambda)c_N/z$.  Although this term defines a perfectly good
classical multiplication operator, directly exponentiating its
quantization would require interpreting a translation that is not small
in $\lambda^{-1}$.  We avoid that interpretation by first using the
normalized characteristic class
\begin{equation}\label{fixed:normalized-class}
 c_0(N_F)=\prod_x(1+x/a)^{-1}.
\end{equation}
Let $Z_{F,\mathrm{ntw}}$ and $\eta_F^{\mathrm{ntw}}$ denote its descendant
potential and pairing.  The subscript distinguishes this theory from
the inverse Euler twist used in localization.

Define $\Delta_{F,0}$ by the formula \eqref{fixed:Delta}, replacing
$-\log(a+x)$ by $-\log(1+x/a)$ and replacing the half logarithm by
$\tfrac12\log(1+x/a)$.  It maps the ordinary paired space to
$(\mathcal H_F,\eta_F^{\mathrm{ntw}})$ and is the identity modulo
$\lambda^{-1}$.

\begin{lemma}[Quantum Riemann--Roch in the normalized twist]
\label{fixed:qrr-normalized}
With the quantization and dilaton translations of
Section~\ref{conv:section},
\begin{equation}\label{fixed:qrr}
 Z_{F,\mathrm{ntw}}=(\text{invertible scalar})\,
                     U_{\Delta_{F,0}}Z_F.
\end{equation}
This is an invertible identity formal in $\lambda^{-1}$ and coefficientwise
in the full Novikov variables and $\hbar$.
\end{lemma}

\begin{proof}
We use the quantum Riemann--Roch theorem of Coates and
Givental~\cite[Theorem~1]{QRR}.  For a multiplicative characteristic class
whose logarithm on a line with first Chern class $x$ is $s(x)$, its
multiplier in the fixed ordinary pairing has exponent
\[
 \sum_x\left\{\frac1z\int_0^x s(t)\,dt+
       \sum_{m\geq1}\frac{B_{2m}}{(2m)!}z^{2m-1}
                                         s^{(2m-1)}(x)\right\}.
\]
The theorem identifies the twisted pairing with the ordinary pairing by
multiplication by $\sqrt{c_0(N_F)}$.  Expressing its conclusion as a map
\emph{to} the actual twisted space therefore adds $-s(x)/2$ to the
exponent.  Taking $s(x)=-\log(1+x/a)$ gives exactly
$\Delta_{F,0}$.

This change of pairing also specifies the affine coordinates of the
quantized action.  In the ordinary-pairing Fock space of the theorem,
the twisted potential is expressed using
\[
 \widetilde q=\sqrt{c_0(N_F)}(t-z1_F).
\]
Multiplication by $c_0(N_F)^{-1/2}$ returns the twisted Darboux
coordinate
\[
 q_{\mathrm{ntw}}=t-z1_F.
\]
Thus the paired-space map
$U_{\Delta_{F,0}}$ has exactly the dilaton translations stipulated in
\eqref{fixed:qrr}.

The theorem applies to the universal-curve complex
$R\pi_*\ev^*N_F$ on the ordinary stable-map spaces of the smooth
projective variety $F$.  Its characteristic class is invertible and
formal in $\lambda^{-1}$.  The quantized multiplier and its inverse are
defined in that filtration, with the dilaton translation in each paired
space.  The scalar factors in quantum Riemann--Roch are immaterial for
projective equations.

The super-space formulation in \cite[Sections~2--3]{QRR} uses the full
cohomology of the target.  In particular, its marked-point terms are
even characteristic-class multiplications and its nodal terms contract
the categorical inverse of the pairing.  It therefore has exactly the
Koszul conventions used here, for arbitrary Hodge types and parity.
\end{proof}

The relation between the two classical multipliers is particularly
simple:
\begin{equation}\label{fixed:Delta-normalization}
 \Delta_F=a^{n/2}
       \exp\left(-\log(a)c_N/z\right)\Delta_{F,0}.
\end{equation}
If $C$ is an ordinary divisor class, then
\[
 [\ell_{0,F},C/z]=[z\partial_z,C/z]+[\mu_F,C/z]=0,
 \qquad [z,C/z]=0.
\]
Here $R_F$ commutes with $C$, and the two displayed contributions to
the first commutator are $-C/z$ and $C/z$.  Consequently $C/z$ commutes
with all the ordinary modes and with their scalar translates
\eqref{fixed:binomial}.  After the coefficient derivatives have been
removed in \eqref{fixed:ordinary-grading}, the scalar $a^{n/2}$ also
cancels in conjugation.  Thus $\Delta_F$ and $\Delta_{F,0}$ produce the
same conjugated mode matrices.  By Lemma~\ref{fixed:qrr-normalized} and
the covariance of Lemma~\ref{conv:covariance}, the equations for those
matrices on $Z_{F,\mathrm{ntw}}$ are equivalent to the ordinary projective
equations on $Z_F$.

It remains to restore the rank factor omitted in
\eqref{fixed:normalized-class}.  This changes the paired space as well
as the invariants, and both changes must be made together.

\begin{lemma}\label{fixed:rank-scaling}
Put $m_F=a^{-n}$.  The actual inverse Euler twist is obtained from the
normalized twist by the substitutions
\begin{equation}\label{fixed:scaling}
 \hbar\longmapsto\hbar/m_F,
 \qquad Q^\beta y^l\longmapsto
     Q^\beta y^l a^{-\int_d c_N},
\end{equation}
where $d$ is the corresponding actual class on $F$.  Its pairing is
$\eta_F^t=m_F\eta_F^{\mathrm{ntw}}$.  Under these substitutions the
quadratic quantizations of
$\Delta_{F,0}\ell_{k,F}\Delta_{F,0}^{-1}$, for every $k\geq-1$, in the
two paired spaces agree.
\end{lemma}

\begin{proof}
On the moduli space of connected genus-$g$ maps of class $d$,
Riemann--Roch gives
\[
 \rank(R\pi_*\ev^*N_F)=n(1-g)+\int_d c_N.
\]
The ratio of the inverse Euler class to the normalized class of this
virtual bundle is therefore
\[
 a^{-n(1-g)-\int_d c_N}
       =m_F^{1-g}a^{-\int_d c_N}.
\]
This is precisely the change in the coefficient of $\hbar^{g-1}Q^\beta
y^l$ under \eqref{fixed:scaling}; exponentiation gives the claimed
identity of total potentials.  The pairing follows by the same
calculation for degree-zero genus-zero three-point invariants.

The $qq$ part of quantization uses the pairing divided by $\hbar$;
the $pp$ part uses its inverse multiplied by $\hbar$.  In the actual
twist these are
\[
 \frac{m_F\eta_F^{\mathrm{ntw}}}{\hbar},\qquad
 \frac{\hbar}{m_F}(\eta_F^{\mathrm{ntw}})^{-1},
\]
which are their normalized counterparts evaluated at $\hbar/m_F$.
The mixed part is unchanged.  The matrices
$\Delta_{F,0}\ell_{k,F}\Delta_{F,0}^{-1}$ have no Novikov derivatives or
Novikov dependence, so the curve substitution commutes with their
action as well.  We use these unshifted matrices here and restore the
scalar $b_F$ afterwards by \eqref{fixed:binomial}.  Every actual curve monomial is multiplied
by a nonzero scalar; hence the substitution is injective and reversible
coefficientwise, without identifying any curve classes.
\end{proof}

The preceding argument used quantum Riemann--Roch only as a formal
identity in $\lambda^{-1}$.  We also need its consequence as an identity
of germs at nonzero $\lambda$.  To make that passage, first remove the
scalar $b_F$ by the invertible binomial change
\eqref{fixed:binomial}.  Each coefficient of the conjugated matrices
$\Delta_{F,0}\ell_{k,F}\Delta_{F,0}^{-1}$ is then rational in $\lambda$.
These matrices have finite lower bounds on their loop powers.  The
coefficients of $Z_{F,\mathrm{tw}}$ are likewise rational in $\lambda$:
the torus acts trivially on $F$, and the inverse Euler class on each
fixed stable-map space is expanded only to its finite cohomological
dimension.  For fixed genus, class, and number of marks, the descendant
indices which can occur are bounded for the same reason.

It follows from the finite-support quantization convention that each
nonconstant coefficient of
\[
 Z_{F,\mathrm{tw}}^{-1}
   \op_{\eta_F^t}
     (\Delta_{F,0}\ell_{k,F}\Delta_{F,0}^{-1})
       Z_{F,\mathrm{tw}}
\]
is a finite sum of rational functions of $\lambda$.  Vanishing of its
formal Laurent expansion at infinity is equivalent to vanishing as a
rational function, and hence as a germ.  At a fixed genus and curve
coefficient, \eqref{fixed:scaling} only shifts finitely many powers of
$\lambda$, so the same argument applies in both directions.  We may now
restore $b_F$ and its logarithms using \eqref{fixed:binomial}.  In
particular, no analytic value of the quantized infinite multiplier is
being taken in this passage.

\begin{proof}[Proof of Proposition~\ref{fixed:equivalence}]
Let $F'$ be the other fixed component.  Up to invertible scalar factors,
the comparison of potentials follows the chain
\[
\begin{gathered}
 Z_F\xrightarrow{\ U_{\Delta_{F,0}}\ }Z_{F,\mathrm{ntw}}
       \xrightarrow{\ \eqref{fixed:scaling}\ }Z_{F,\mathrm{tw}},
 \\
 Z_{F,\mathrm{tw}}\xrightarrow{\ U_{S_f}\ }
       \mathcal A_{F,\mathrm{tw}}(u_F),
 \\
 \mathcal A_{F,\mathrm{tw}}(u_F)
 \mathcal A_{F',\mathrm{tw}}(u_{F'})
       \xrightarrow{\ U_R\ }\mathcal A_W.
\end{gathered}
\]
The arrow labeled \eqref{fixed:scaling} changes parameters and rescales
the pairing as in Lemma~\ref{fixed:rank-scaling}.  The final row first
adjoins the other fixed potential and then applies the upper symplectic
transformation.

Lemma~\ref{fixed:classical}, the binomial identity, quantum
Riemann--Roch, and Lemma~\ref{fixed:rank-scaling} identify projective
satisfaction of all ordinary modes on $Z_F$ with projective
satisfaction of the corresponding descendant modes on
$Z_{F,\mathrm{tw}}$.  The preceding rationality argument gives this
equivalence over the same germs in $\lambda$ used by the spectral
construction.

Pass next from descendants on $F$ to ancestors at $u_F$.  The
ancestor--descendant identity of Section~\ref{geo:ancestors} and
Lemma~\ref{conv:covariance} conjugate the mode matrices by
\[
 S_f=S_{F,\mathrm{tw}}(u_F,z)
\]
and transport their projective equations to
$\mathcal A_{F,\mathrm{tw}}(u_F)$.  On the product of the two fixed
potentials, an operator supported on $F$ acts only on that factor; its
projective equation is consequently equivalent to the equation on
the single factor.

Finally apply the upper transformation $R(z)$ of
Proposition~\ref{loc:factorization}.  Its two identities
\[
 S_W=R(z)\bigoplus_F S_f,
 \qquad
 \mathcal A_W=(\text{scalar})\,
      U_R\prod_F\mathcal A_{F,\mathrm{tw}}(u_F)
\]
show that the resulting classical modes are exactly $A_{k,F}$, by
Proposition~\ref{spec:zero}, and that their quantum equations are the
projective equations on $\mathcal A_W$.  Each change is invertible, so
the implication holds in both directions.  The classical changes of
paired space are symplectic, proving also the symplectic assertion in
the proposition.

All the quantized changes in this last paragraph are defined in the
Novikov completion: $u_F$, $S_f-I$, and $R-I$ have positive Novikov
filtration.  The fixed-theory series $S_f$ has a finite negative tail
at each curve coefficient, and $R$ is upper.  Conjugation and
quantization therefore preserve the lower loop bounds and the finite
support required in Section~\ref{conv:section}.  In particular, the
comparison uses the upper quantization supplied by localization and
the fixed-theory ancestor changes; it does not require quantizing a
new expansion of $S_W$ itself.
\end{proof}
